%% pengesahan.tex
%% Lembar Pengesahan Tesis
%% Format sesuai template ITS dengan background pattern dan logo watermark
%% Font: Times New Roman 12pt
%%
%% EDIT DATA DI BAWAH INI SESUAI DATA TESIS ANDA

\newpage
\thispagestyle{plain}

% Background and watermark using TikZ
\begin{tikzpicture}[remember picture,overlay]
    % Background pattern - full page
    \node[anchor=center, inner sep=0pt] at (current page.center) {
        \includegraphics[width=\paperwidth,height=\paperheight]{figures/background-pengesahan.png}
    };
    
    % ITS Logo watermark - centered, bigger (logo sudah transparan)
    \node[anchor=center] at (current page.center) {
        \includegraphics[width=13cm]{figures/logo-pengesahan.png}
    };
\end{tikzpicture}

\vspace*{-0.5cm}

% Content - reduced spacing to fit 1 page
\begin{center}

%% ===== JUDUL =====
{\fontsize{14pt}{16pt}\selectfont\bfseries LEMBAR PENGESAHAN TESIS}

\vspace{0.3cm}

%% ===== KETERANGAN =====
{\fontsize{12pt}{14pt}\selectfont
Tesis disusun untuk memenuhi salah satu syarat memperoleh gelar,\\
Magister Komputer (M.Kom.)\\
di\\
\textbf{Institut Teknologi Sepuluh Nopember}}

\vspace{0.3cm}

%% ===== DATA MAHASISWA =====
{\fontsize{12pt}{14pt}\selectfont
\textbf{Oleh:}\\
\colorbox{yellow}{Nama Mahasiswa}\\
NRP. \colorbox{yellow}{6025XXXXXXX}}

\vspace{0.2cm}

%% ===== TANGGAL DAN PERIODE =====
{\fontsize{12pt}{14pt}\selectfont
Tanggal Ujian: \colorbox{yellow}{dd mm yyyy}\\
Periode Wisuda: \colorbox{yellow}{mm xxxx}}

\vspace{0.2cm}

%% ===== DISETUJUI =====
{\fontsize{12pt}{14pt}\selectfont\textbf{Disetujui oleh:}}

\end{center}

\vspace{0.1cm}

%% ===== DAFTAR PEMBIMBING DAN PENGUJI =====
{\fontsize{12pt}{14pt}\selectfont
\begin{enumerate}[itemsep=0.08cm, topsep=0pt, parsep=0pt]
    \item \begin{tabular}[t]{@{}p{8.5cm}p{4.5cm}@{}}
        Nama Pembimbing 1 & \dotfill \\[-0.15cm]
        NIP. XXXXX & \\[-0.15cm]
        (Pembimbing 1) &
    \end{tabular}
    
    \item \begin{tabular}[t]{@{}p{8.5cm}p{4.5cm}@{}}
        Nama Penguji 1 & \dotfill \\[-0.15cm]
        NIP. XXXXX & \\[-0.15cm]
        (Penguji 1) &
    \end{tabular}
    
    \item \begin{tabular}[t]{@{}p{8.5cm}p{4.5cm}@{}}
        Nama Penguji 2 & \dotfill \\[-0.15cm]
        NIP. XXXXX & \\[-0.15cm]
        (Penguji 2) &
    \end{tabular}
    
    \item \begin{tabular}[t]{@{}p{8.5cm}p{4.5cm}@{}}
        Nama Penguji 3 & \dotfill \\[-0.15cm]
        NIP. XXXXX & \\[-0.15cm]
        (Penguji 3) &
    \end{tabular}
\end{enumerate}
}

\vspace{0.3cm}

%% ===== KEPALA DEPARTEMEN =====
%% Rata kiri sejajar dengan nomor list menggunakan hspace di setiap baris
{\fontsize{12pt}{14pt}\selectfont
\noindent\hspace{0.5cm}Kepala Departemen Teknik Informatika

\vspace{0.6cm}

\noindent\hspace{0.5cm}Nama Kepala Departemen

\noindent\hspace{0.5cm}NIP. XXXXX
}